\originalchapter{相互作用玻色子与相位动力学}
\label{ch:boson}
\source{8.513第7--15讲组：相互作用玻色子、超流与低维有效理论。}
\originalsection{相干态与路径积分}
玻色相干态满足$b\ket z=z\ket z$，完备关系$\int\dd^2z\,\ket z\bra z/\pi=I$（采用归一化态及对应权重）。把$\e^{-\beta H}$分为很多小虚时间片并逐片插入相干态，相邻态重叠产生一阶时间导数。对正规序哈密顿量得到
\begin{equation}
Z=\int\mathcal D(z^*,z)\e^{-S_E/\hbar},\qquad S_E=\int_0^{\beta\hbar}\dd\tau\,[\hbar z^*\partial_\tau z+H(z^*,z)].
\end{equation}
玻色场虚时间周期；费米相干态使用Grassmann变量且反周期，不能用普通复变量的积分规则替代。经典实时间作用量为$L=\ii\hbar z^*\dot z-H$，差一个总导数可把Berry项写为$\hbar n\dot\theta$。本文取$z=\sqrt n\e^{-\ii\theta}$以使该项为正，若改用$\e^{\ii\theta}$，相位定义和共轭动量应同时改号。
\originalsection{Bogoliubov声子}
均匀弱排斥玻色气$H=\int\dd^dx[\hbar^2|\nabla\psi|^2/(2m)+g|\psi|^4/2-\mu|\psi|^2]$。均匀鞍点$n_0=\mu/g$，写$\psi=\sqrt{n_0}+\delta\psi$并保留二次项，$k,-k$块含$(\epsilon_k+gn_0)b_k^\dagger b_k$与$gn_0(b_kb_{-k}+\mathrm{h.c.})/2$。保持对易关系的变换$u_k^2-v_k^2=1$给
\begin{equation}
E_k=\sqrt{\epsilon_k(\epsilon_k+2gn_0)},\quad u_k^2,v_k^2=\frac12\left(\frac{\epsilon_k+gn_0}{E_k}\pm1\right),\quad\epsilon_k=\frac{\hbar^2k^2}{2m}.
\end{equation}
低$k$线性$E_k\simeq\hbar c|k|$，$c=\sqrt{gn_0/m}$；高$k$回到粒子色散。声子是凝聚密度与相位耦合的集体模。$v_k^2$给零温耗尽，计算时若接触势在高维产生紫外问题，应以散射长度重整化$g$，不能无限外推裸接触模型。
\originalsection{密度--相位有效作用量}
写$\psi=\sqrt{n_0+\delta n}\e^{-\ii\theta}$，低梯度保留$L=\int\dd x[\hbar(n_0+\delta n)\dot\theta-g(\delta n)^2/2-\hbar^2n_0(\nabla\theta)^2/(2m)]$。密度涨落无时间导数，是相位的共轭变量。极小化$\delta n=\hbar\dot\theta/g$并代回，得到
\begin{equation}
L_{\rm eff}=\int\dd x\left[\hbar n_0\dot\theta+\frac{\hbar^2}{2g}\dot\theta^2-\frac{\hbar^2n_0}{2m}(\nabla\theta)^2\right].
\end{equation}
第一项在光滑小涨落中是总导数，但对紧致相位的绕数与扭曲边界有作用，不能随意删去。Euler--Lagrange方程给$\ddot\theta-c^2\nabla^2\theta=0$；声速与Bogoliubov结果相同，证明两个描述在低能一致。

一维常写$S_E/\hbar=(K/2\pi)\int\dd x\dd\tau[v^{-1}(\partial_\tau\theta)^2+v(\partial_x\theta)^2]$。高斯反演在二维$(x,v\tau)$给$\avg{[\theta(x)-\theta(0)]^2}=K^{-1}\log(|x|/a)$，于是$\avg{\e^{\ii\theta(x)}\e^{-\ii\theta(0)}}\propto|x|^{-1/(2K)}$。无穷系统的相位方差发散，所以零温一维并非具有常数远距关联的通常凝聚，而是代数序。有限温虚时间有有限周长，长距离关联进一步衰减。
\originalsection{零模、扭曲与LSM约束}
环长$L$的相位写$\theta=2\pi mx/L+\theta_0(t)+\theta_{\ne0}$。取$L_{\rm eff}=\int[A\dot\theta^2/2-B(\partial_x\theta)^2/2+n_0\hbar\dot\theta]$，零模角动量$p_0=AL\dot\theta_0+n_0\hbar L$，因$\theta_0$为$2\pi$周期坐标而$p_0=\hbar N$，$N\in\mathbb Z$。零模能量
\begin{equation}
E_{N,m}=\frac{\hbar^2(N-n_0L)^2}{2AL}+\frac{2\pi^2Bm^2}{L},\qquad P_{N,m}=\frac{2\pi\hbar mN}{L}.
\end{equation}
末式由平移引起相位整体移$2\pi m\Delta x/L$获得。密度Berry项虽不改声子方程，却改变量子数与扭曲后的动量。

格点环上定义$U=\exp[(2\pi\ii/L)\sum_jj n_j]$，整数局域电荷使边界因子为1，平移关系给$TU T^{-1}=U\exp[-2\pi\ii N/L]$（若平移方向反定，符号反转）。故$U\ket0$的动量与基态差$2\pi\nu$模$2\pi$。局域跳跃经$U$只获得$O(1/L)$相位；正负扭曲合并可消去一阶电流项，变分能量增加至多$O(1/L)$。若$\nu=p/q$既约且$q>1$，重复扭曲产生$q$个不同动量的低能态，排除在保持平移与$U(1)$的条件下唯一且有统一能隙的基态。结论可能由简并、破缺或无隙实现，不能把它只说成“必然超流”。该推导要求局域性、适当的局域能量界和守恒电荷，非局域作用不自动满足。
